% !PRINTED 779
\weberpage{779}
% running head "Primzahlen in Linearformen." / "779" at the top of this page is not
% transcribed as body text.

Now
\[
\frac{\mu_x}{m} - x = \frac{a-m}{m},
\]
and consequently we can apply the theorem \S.\,191, 6. to the quantities
$\mu_x$, in which we have to set $n = x$, $\gamma = 1:m$, $c_n = (a-m):m$.
Accordingly we have
\weq{7}{A(s) = \overset{x}{\underset{0,\,\infty}{\Sigma}} \frac{1}{(mx+a)^s}
= \overset{x}{\underset{1,\,\infty}{\Sigma}} \frac{1}{(mx+a-m)^s},}
which, so long as $s > 1$, is a continuous function of $s$, and we have
\weq{8}{A(s) = \frac{1}{m(s-1)} + C(s),}
where $C(s)$ remains finite for $s=1$\footnote{By the use of known
theorems on $\Gamma$-functions one finds, from \S.\,191, (20):
\[
C(1) = \frac{1}{m} \frac{\Gamma'\left(\frac{a}{m}\right)}{\Gamma\left(\frac{a}{m}\right)}.
\]}.

We now let $A$ run through all the classes (2), denote by
$\chi$ any one of the characters of the group $\mathfrak{N}$, and set
\weq{9}{Q(s) = \overset{A}{\Sigma} \chi(A)\, A(s),}
or, what is the same thing,
\weq{10}{Q(s) = \overset{n}{\Sigma} \frac{\chi(n)}{n^s},}
in which the sum extends over all the numbers $n$. The number of these sums
$Q$ is as great as the number of the characters, that is, $= \varphi(m)$. We
denote them, when a distinction is necessary, by
\[
Q_1, Q_2, \ldots, Q_\mu.
\]

One of them, $Q_1$, corresponds to the principal character and has the expression
\[
Q_1(s) = \overset{n}{\Sigma} \frac{1}{n^s}.
\]

Now [\S.\,11, (21)]:
\[
\overset{A}{\Sigma} \chi(A) = \mu \ \text{or} = 0,
\]
according as $\chi$ is the principal character or not. Furthermore, there
results from (8)
%CONT>>
